\PassOptionsToPackage{hyphens}{url}
\documentclass[final,5p,twocolumn,a4paper]{elsarticle}

\usepackage{amssymb}
\usepackage{amsmath}
\usepackage{hyperref}
\usepackage{graphicx}
\usepackage{booktabs}
\usepackage{listings}
\usepackage{xcolor}
\usepackage{float}
\usepackage{url}

\setlength{\textfloatsep}{8pt plus 2pt minus 2pt}
\setlength{\intextsep}{8pt plus 2pt minus 2pt}
\setlength{\dbltextfloatsep}{8pt plus 2pt minus 2pt}
\setlength{\floatsep}{8pt plus 2pt minus 2pt}
\setlength{\dblfloatsep}{8pt plus 2pt minus 2pt}

\definecolor{codegreen}{rgb}{0,0.6,0}
\definecolor{codegray}{rgb}{0.5,0.5,0.5}
\definecolor{codepurple}{rgb}{0.58,0,0.82}
\definecolor{backcolour}{rgb}{0.96,0.96,0.96}

\lstdefinestyle{mystyle}{
    backgroundcolor=\color{backcolour},   
    commentstyle=\color{codegreen},
    keywordstyle=\color{blue},
    numberstyle=\tiny\color{codegray},
    stringstyle=\color{codepurple},
    basicstyle=\ttfamily\scriptsize,
    breaklines=true,                 
    captionpos=b,                    
    keepspaces=true,                 
    numbers=left,                    
    numbersep=5pt,                  
    showspaces=false,                
    showstringspaces=false,
    showtabs=false,                  
    tabsize=2
}
\lstset{style=mystyle}

\journal{SoftwareX}

\begin{document}
\renewcommand{\labelenumii}{\arabic{enumi}.\arabic{enumii}}

\begin{frontmatter}

\title{LunarAtlas: A Reproducible Data Processing and Interactive Visualization Infrastructure for Chandrayaan-3 LIBS Level-1 Archives}

\author{Lovekesh Anand\corref{cor1}}
\ead{lovekeshanand6@gmail.com}
\author{Dua Saeed}
\ead{23f1000825@ds.study.iitm.ac.in}

\address{Department of Computer Science and Engineering, Maharshi Dayanand University, Rohtak, Haryana, India}
\cortext[cor1]{Corresponding author}

\begin{abstract}
LunarAtlas is an open-source, reproducible data processing infrastructure that transforms publicly available Chandrayaan-3 laser-induced breakdown spectroscopy (LIBS) Level-1 (L1) products from the ISRO Planetary Data Archive into normalised, analysis-ready spectral records suitable for quantitative lunar science. The Chandrayaan-3 LE-LIBS L1 archive distributes calibrated spectra in wide-format tables where 2,094 wavelength channels (instrument range: 220--800~nm; stored channel range: 164.35--878.26~nm) appear as individual columns, with plasma and background acquisitions interleaved. To our knowledge, no publicly documented, reproducible pipeline exists to ingest and clean these products automatically. LunarAtlas implements a six-stage deterministic Python pipeline that: (i) parses PDS4 XML metadata; (ii) verifies file provenance via MD5 checksums; (iii) reshapes wide-format records into a normalised long-form relational layout $(\text{Measurement\_ID},\,\text{Wavelength\_nm}) \rightarrow \text{Intensity (a.u.)}$; (iv) identifies plasma and background shots from mission status flags; (v) applies paired background subtraction with clamping; and (vi) assigns file-unique Measurement IDs. For interactive visualisation, LunarAtlas employs an LTTB+Peaks adaptive downsampling algorithm that retains 100\% of detected spectral peaks. A versioned REST API delivers spectral data; under single-connection, no-load conditions on a commodity workstation, the measured API response latency was below 500~ms --- a hardware-specific benchmark result that may vary under concurrent or networked deployment. Software and documentation are available at: \url{https://github.com/LovekeshAnand/LunarAtlas} (archived snapshot: \url{https://doi.org/10.5281/zenodo.20771541}).
\end{abstract}

\begin{keyword}
LIBS \sep Chandrayaan-3 \sep planetary spectroscopy \sep PDS4 \sep database \sep downsampling \sep LTTB \sep REST API
\end{keyword}

\end{frontmatter}

\section*{Required Metadata}

\begin{table*}[!ht]
\centering
\small
\begin{tabular}{|l|p{7.5cm}|p{7.5cm}|}
\hline
\textbf{Nr.} & \textbf{Code metadata description} & \textbf{Please fill in this column} \\
\hline
C1 & Current code version & v1.0.0 \\
\hline
C2 & Permanent GitHub link to code/repository used for this code version & \href{https://github.com/LovekeshAnand/LunarAtlas}{https://github.com/LovekeshAnand/LunarAtlas} \\
\hline
C3 & Legal Code License & MIT License \\
\hline
C4 & Code versioning system used & git \\
\hline
C5 & Software code languages, tools, and services used & Python 3.11, SQL, TypeScript, HTML, CSS \\
\hline
C6 & Compilation requirements, operating environments \& dependencies & Python~3.11, PostgreSQL~16.3, Node.js~18.x, Redis; dependencies: numpy, pandas, scipy, fastapi, React \\
\hline
C7 & If available Link to developer documentation/manual & \href{https://github.com/LovekeshAnand/LunarAtlas}{https://github.com/LovekeshAnand/LunarAtlas} \\
\hline
C8 & Support email for questions & lovekeshanand6@gmail.com, 23f1000825@ds.study.iitm.ac.in \\
\hline
\end{tabular}
\caption{Code metadata (mandatory)}
\label{table:metadata} 
\end{table*}

\section{Motivation and Significance}

Laser-Induced Breakdown Spectroscopy (LIBS) is a well-established technique for in-situ elemental analysis in planetary exploration. By focusing a pulsed laser onto a target surface, LIBS generates a transient plasma whose emission spectrum reveals the elemental composition of the sample. The technique has been deployed on multiple planetary missions, including ChemCam on MSL \cite{maurice2012chemcam, wiens2012chemcam}, SuperCam on Mars 2020 \cite{maurice2021supercam, wiens2021supercam}, and MarSCoDe on Tianwen-1 \cite{xu2021marscode}. In 2023, India's Chandrayaan-3 mission operated the Laser-Induced Breakdown Spectroscopy (LIBS) payload aboard the Pragyan rover, acquiring the first in-situ LIBS spectra from the lunar south polar region \cite{megala2024chandrayaan, sridhar2024libs, isro_chandrayaan3_libs}.

The Pragyan LIBS instrument employs a pulsed Yb:Er:Glass laser operating at $1535 \pm 1$~nm and records plasma emission spectra over the calibrated 220--800~nm spectral range using a holographic flat-field grating spectrometer coupled to a linear CCD detector. The archived Level-1 (L1) products contain 2,094 wavelength channels spanning 164.35--878.26~nm, extending beyond the instrument's calibrated operating range. Throughout this work, all scientific interpretation is restricted to the calibrated spectral interval. The instrument provides a spectral sampling of approximately 0.34~nm per channel with an optical resolution of $\le 1$~nm. These datasets are publicly distributed through the ISRO Planetary Data Archive (PRADAN) as Planetary Data System Version 4 (PDS4) archives \cite{isro_chandrayaan3_libs, pds4_standard}.

Although the archive ensures long-term preservation and public accessibility, the delivered L1 products are not directly suitable for reproducible computational analysis. Several characteristics of the archive require mission-specific preprocessing before the spectra can be queried, analysed, or visualised efficiently.
\begin{enumerate}
    \item \emph{Interactive Visualisation:} Each spectrum contains 2,094 sampled wavelength channels. Efficient web-based visualisation therefore benefits from adaptive downsampling techniques that preserve diagnostically important emission peaks while reducing rendering overhead.
    \item \emph{Tabular Reshaping:} Each spectrum is distributed as a wide-format table in which all 2,094 wavelength channels are stored as individual columns. This representation complicates relational database normalisation, prevents efficient wavelength-based queries, and limits scalable downstream analysis.
    \item \emph{Interleaved Acquisitions:} Plasma spectra and corresponding dark-field background measurements are interleaved within each L1 product according to the PRADAN acquisition sequence. Recovering cleaned plasma spectra requires parsing mission status flags and correctly pairing each plasma acquisition with its corresponding background measurement.
    \item \emph{Background Suppression:} Quantitative analysis requires subtracting the paired background spectrum from each plasma spectrum. Negative intensities produced after subtraction are clamped to zero to suppress non-physical values arising from detector noise and measurement uncertainty.
\end{enumerate}

To our knowledge, no publicly documented open-source software currently provides an end-to-end processing workflow specifically for Chandrayaan-3 LIBS Level-1 products. Table~\ref{table:capability_comparison} compares LunarAtlas with existing general-purpose and related software packages: \emph{pds4\_tools} \cite{pds4_tools_ref} and \emph{libpyhat} \cite{anderson2019libpyhat}.

\begin{table*}[!htbp]
\centering
\small
\setlength{\tabcolsep}{6pt}
\begin{tabular}{lccc}
\toprule
\textbf{Capability} & \textbf{pds4\_tools (v1.4)} \cite{pds4_tools_ref} & \textbf{libpyhat (v0.x)} \cite{anderson2019libpyhat} & \textbf{LunarAtlas (v1.0.0)} \\
\midrule
PDS4 XML label parsing & Yes & No & Yes \\
Calibrated spectral axis reconstruction & No & No & Yes \\
Plasma/background flag parsing ($F$, $L$ flags) & No & No & Yes \\
Paired background subtraction & No & No & Yes \\
Unmatched plasma shot handling (file boundary) & No & No & Yes \\
Wide-to-long relational reshaping & No & No & Yes \\
MD5 provenance verification & No & No & Yes \\
Spectral preprocessing (norm, baseline, PCA) & No & Yes & Partial$^\dagger$ \\
Relational database ingestion & No & No & Yes \\
REST API + Interactive visualisation & No & No & Yes \\
\bottomrule
\multicolumn{4}{p{0.95\textwidth}}{\footnotesize $^\dagger$ LunarAtlas exposes ALS and Savitzky-Golay denoising as optional post-ingestion API filters; normalisation, PCA, and machine-learning workflows remain the domain of tools such as libpyhat, with which LunarAtlas output is designed to be compatible via CSV export.}
\end{tabular}
\caption{Capability comparison between general-purpose PDS4 reading tools, general-purpose spectral analysis frameworks, and LunarAtlas.}
\label{table:capability_comparison}
\end{table*}

While \emph{pds4\_tools} successfully parses the PDS4 XML label and locates the associated CSV via byte-offset references, it delivers the spectral payload as an unstructured delimited table with no wavelength-to-intensity axis construction. The 2,094 \texttt{Response\_Count} repetitions defined in the XML label carry no wavelength calibration; those values exist only as column headers in the CSV itself. Consequently, \emph{pds4\_tools} provides no facility for reconstructing the calibrated spectral axis, structuring the eight interleaved acquisitions into a (record $\times$ wavelength) matrix, or applying any LIBS-specific preprocessing steps.

Similarly, while \emph{libpyhat} provides documented capabilities for spectral preprocessing, normalisation, baseline removal, dimensionality reduction, and machine-learning workflows, its ingestion architecture is structurally incompatible with Chandrayaan-3 LE-LIBS L1 products. The \texttt{SpectralData} constructor requires input as a pandas DataFrame carrying a two-level MultiIndex with mandatory top-level labels \texttt{wvl}, \texttt{meta}, or \texttt{comp}. The raw L1 CSV carries a flat single-row header of the form \texttt{['Time', 'Count', 'Reset', 'Fire', '164.35', ..., '878.26']}. Because no column in this header is a tuple, the constructor's format-gate drops every column and raises an error before any processing method can be reached, resulting in a hard crash at object instantiation rather than a graceful degradation. Beyond this format incompatibility, \emph{libpyhat} provides no facility for parsing PDS4 XML labels, interpreting Force Reset Status ($F$) and Laser Fire Status ($L$) flags to distinguish plasma from background acquisitions, or pairing and subtracting interleaved shots. Thus, \emph{libpyhat} operates as a processing and analysis engine for spectra that have already been ingested and cleaned; LunarAtlas fills the upstream gap between raw ISRO PDS4 delivery and the point where analytical toolkits can be applied.

LunarAtlas addresses this software gap by providing an open-source, deterministic, and reproducible processing infrastructure that automates PDS4 metadata parsing, archive integrity verification, plasma--background pairing, background subtraction, relational database ingestion, REST API generation, and interactive spectral visualisation. By transforming archived Chandrayaan-3 LIBS products into analysis-ready spectral records, LunarAtlas enables reproducible downstream scientific workflows while eliminating the need for custom mission-specific preprocessing software.

\section{Software description}

\subsection{Software Architecture}

LunarAtlas adopts a three-tier architecture. The overall system flow is shown in Figure~\ref{fig:architecture}. The layers consist of:
\begin{itemize}
    \item \textbf{Data Layer:} A PostgreSQL database \cite{postgresql} stores mission metadata, observation provenance, and processed spectra. Processed values are stored in a normalised long-form layout, where each record represents a single channel value.
    \item \textbf{Logic Layer:} Implemented in Python~3.11 and FastAPI \cite{fastapi}, the logic layer manages NIST atomic spectra matching, adaptive downsampling, and API query routing. All numerical reductions---including spectral normalisation, dark subtraction, and peak integration---are executed server-side in Python to prevent floating-point divergence that would arise from equivalent operations in JavaScript. Result caching is handled via Redis \cite{redis}, which stores recently served downsampled payloads using a Least Recently Used (LRU) eviction policy with a 1-hour Time-To-Live (TTL); keys expire automatically after this window, eliminating the need for manual cache flushing. Since the underlying Level-1 observation records are static, cache invalidation across software versions is managed via deterministically constructed version-prefixed keys of the form \texttt{lunaratlas\_v\{n\}:\{endpoint\}:\{parameter\_hash\}}, where \texttt{n} is the software version integer and \texttt{parameter\_hash} is an MD5 digest of the query parameters. Bumping \texttt{n} upon deployment renders all prior-version keys unreachable without affecting other Redis tenants on the same instance; orphaned keys from retired versions are removed on a scheduled basis using \texttt{redis-cli --scan --pattern "lunaratlas\_v\{n\}:*"} | \texttt{xargs redis-cli del}.
    \item \textbf{Presentation Layer:} A browser-based React client \cite{react} renders interactive spectral plots. To ensure scientific integrity, all data reduction (including background subtraction, non-negativity clamping, and LTTB+Peaks downsampling) remains server-side. This architecture prevents floating-point arithmetic divergence between browser-based JavaScript engines and Python-based scientific libraries, ensuring that the visualised spectra match the ingested database records bit-for-bit.
\end{itemize}



Dependency versions are pinned in \texttt{requirements.txt} (Python~3.11) and \texttt{package.json} (Node.js~18). A multi-stage Docker image is provided at the repository root (\texttt{Dockerfile}), which builds the Vite frontend and the FastAPI backend together for simplified single-container deployment. Separate \texttt{Dockerfile.backend} and \texttt{Dockerfile.frontend} images are also provided for independent service deployment. A tagged release snapshot of LunarAtlas~v1.0.0 is archived at Zenodo~\cite{lunaratlas_zenodo} to satisfy reproducibility requirements. To ensure computational reproducibility and allow offline verification without PRADAN archive access, a comprehensive unit test suite is included in the repository under the \texttt{tests/} directory. Running \texttt{pytest tests/ -v} executes automated smoke and unit tests against synthetic input signals (e.g., simulated plasma/background spectrum pairs and baseline noise). These tests validate the core pipeline mathematics, including non-negativity clamping, paired background subtraction, LTTB downsampling ratios, and NIST peak-lock union logic.

\begin{figure*}[!t]
    \centering
    \includegraphics[width=0.9\textwidth]{publication_figures/fig1_architecture}
    \caption{End-to-end LunarAtlas system architecture and data flow schematic.}
    \label{fig:architecture}
\end{figure*}

As illustrated in Figure~\ref{fig:architecture}, the data ingestion and query flow proceeds sequentially: raw Level-1 PDS4 data are processed offline and stored in PostgreSQL (Data Layer); a FastAPI backend serves these records, applying optional denoising and downsampling under API key authentication and rate-limiting (Logic Layer); and a React client renders interactive plots in the browser (Presentation Layer).


\subsection{Data Processing and Cleaning Pipeline}

The processing pipeline consists of eight sequential scripts: folder structure validation, L1 data processing (the six internal stages described below), QA plot generation, NIST peak verification, MD5 signing of processed outputs, data folder segregation, database ingestion, and a final verification audit. The processing pipeline is implemented in Python across two primary modules: \texttt{step2\_process\_l1\_data.py} (ingestion, pairing, cleaning) and \texttt{step3\_graph\_plotting.py} (QA plot generation). The pipeline is fully deterministic and executes in six sequential stages to transform raw PDS4 Level-1 (L1) products into cleaned relational records:

\begin{enumerate}
    \item \textbf{PDS4 Label Parsing:} The pipeline parses the accompanying PDS4 XML label to extract logical identifiers, observation start/stop times, physical units, and instrument metadata. The parser follows PDS4 field semantics \cite{pds4_standard} to locate mandatory fields including \texttt{pds:logical\_identifier}, \texttt{pds:start\_date\_time}, and \texttt{pds:stop\_date\_time}; labels from which these fields cannot be extracted are skipped with a warning.
    \item \textbf{File Provenance Verification:} A 128-bit MD5 checksum is computed over the raw L1 CSV bytes before any processing begins. MD5 is used for file integrity verification in a trusted archive context, not for cryptographic security. The digest is stored alongside ingestion timestamps and algorithm version strings, establishing a traceable audit trail. The pipeline is fully deterministic: re-running on the same input file produces bit-identical output records. A corresponding MD5 manifest is generated for all processed output files (Stage 5) and cross-verified against database records in a final audit stage (Stage 8), establishing end-to-end provenance across the full pipeline.
    \item \textbf{Wide-to-Long Reshaping:} Spectra are distributed as wide CSV tables in which each wavelength channel (164.35--878.26~nm; 2,094 channels total) is represented as a separate column \cite{numpy, pandas}. The pipeline reshapes each row into a normalised long-form layout, as defined by the relation in Equation~\eqref{eq:reshape}:
    \begin{equation}
        \label{eq:reshape}
        (\text{Measurement\_ID}, \text{Wavelength\_nm}) \rightarrow \text{Intensity (a.u.)}
    \end{equation}
    This relational representation eliminates position-based column assumptions and supports efficient range-based queries. Prior to reshaping, the pipeline validates the presence of eight trailing instrument parameter columns (Delay Time, Integration Time, Number of Pulses, X Factor, Laser Energy, Laser Pump Current, PRR, On Off Status); files failing this check are skipped with an error.


    \item \textbf{Plasma-Background Identification:} The pipeline parses Force Reset Status ($F$) and Laser Fire Status ($L$) metadata flags to identify background ($B$) and plasma ($P$) shots according to the boolean logic in Equation~\eqref{eq:flags}:
    \begin{equation}
        \label{eq:flags}
        B: (F = 1) \land (L = 0) \quad \text{and} \quad P: (F = 0) \land (L = 1)
    \end{equation}
    Background association relies on acquisition order preserved in the L1 CSV. Each plasma acquisition is associated with the most recently observed valid background acquisition, which is ordered per the PRADAN file specification. Acquisitions satisfying neither condition are considered anomalous and excluded. When consecutive plasma shots appear without an intervening background, each is subtracted from the most recently seen background acquisition. The pipeline additionally handles: empty files; background-only or plasma-only files; ambiguous status flag combinations (0,0) or (1,1); NaN values in status or intensity columns; missing XML label files; and trailing background acquisitions with no subsequent plasma shot — all logged with provenance warnings and excluded from output.
    \item \textbf{Background Subtraction and Clamping:} Each plasma shot $P_i$ is paired with its most recently preceding background shot $B_i$, under the assumption of measurement stationarity within a single acquisition cycle. This assumption is justified by the Pragyan LIBS instrument's fixed acquisition cadence of 5 Hz (200 ms inter-shot interval within a single command sequence), as specified in the instrument technical documentation \cite{libs_data_manual, libs_sis_ref}, during which ambient temperature, detector dark current, and surface characteristics remain constant. The clean spectrum is computed channel-wise according to Equation~\eqref{eq:clamping}:
    \begin{equation}
        \label{eq:clamping}
        I_{clean}(\lambda) = \max(0, I_{P_i}(\lambda) - I_{B_i}(\lambda))
    \end{equation}
    Intensity values represent detector response arbitrary units (a.u.) and must be non-negative; negative values arising from noise fluctuations are set to exactly zero. Acquisitions in which more than 50\% of channels are clamped to zero are flagged with a quality warning (e.g., Measurement ID 2, 77.4\% clamped channels, Table 4). NaN intensity values in plasma spectra, arising from detector readout failures, are replaced with 0.0 prior to subtraction; affected acquisitions are flagged via the Has\_NaN\_Intensities field in the output record.
    \item \textbf{Measurement ID Generation:} To resolve timestamp collisions (when multiple acquisitions share a single clock second), the pipeline assigns a sequential, file-unique Measurement ID. IDs are unique within a file; during database ingestion, a globally unique identifier is constructed by combining the observation-specific file identifier with the local Measurement ID, enabling unambiguous retrieval of spectra across the entire archive. This ID groups the 2,094 wavelength-intensity records into a single logical entity. The globally unique identifier takes the form \{file\_info\_id\}-\{measurement\_index\}, where file\_info\_id is the database primary key of the parent observation file record.
\end{enumerate}




\subsubsection{Justification of Paired Background Subtraction}
Paired subtraction was selected over model-based baseline methods for three reasons: (i) the LE-LIBS acquisition protocol provides temporally matched background measurements, making paired subtraction directly applicable without additional modelling assumptions \cite{sridhar2024libs}; (ii) the operation is fully deterministic and parameter-free, satisfying the reproducibility requirements of the infrastructure layer; and (iii) downstream users may apply further denoising (e.g., ALS \cite{eilers2003perfect} or Savitzky-Golay \cite{savitzky1964}) without the cleaning tier imposing irreversible smoothing choices. Table~\ref{table:subtraction_tradeoffs} summarises the trade-offs.

\begin{table*}[!htbp]
\centering
\footnotesize
\setlength{\tabcolsep}{4pt}
\begin{tabular}{@{}lp{1.6cm}cc@{}}
\toprule
\textbf{Method} & \textbf{Tuning params} & \textbf{Determ.} & \textbf{Acq. aware} \\
\midrule
Paired subtraction (LunarAtlas) & None & Yes & Yes \\
ALS \cite{eilers2003perfect} & $\lambda$, p & No & No \\
SWT \cite{nason1995stationary} & level, basis & No & No \\
Savitzky--Golay \cite{savitzky1964} & window, order & Yes & No \\
\bottomrule
\end{tabular}
\caption{Comparison of background-removal strategies for LE-LIBS L1 data.}
\label{table:subtraction_tradeoffs}
\end{table*}

\subsubsection{LTTB+Peaks Adaptive Downsampling}
A full-resolution LE-LIBS spectrum contains 2,094 channels. Rendering thousands of spectra simultaneously in a browser-based visualisation console creates significant client-side overhead. Downsampling is therefore required for interactive visualisation while preserving diagnostically important spectral features.

LunarAtlas downsamples the data to a target count $N(k) = \min(2094, \lceil 2094 \cdot 2^{-k} \rceil)$, where $k$ is the zoom level ($k=0$ returns the full 2,094 points). The Largest Triangle Three Buckets (LTTB) algorithm \cite{steinarsson2013lttb} partitions the spectrum into buckets and selects points that maximise the triangle area formed with adjacent buckets. For a candidate point $P_c$ under evaluation in bucket $j$, LTTB selects the point that maximises the area of the triangle formed with two anchor points: $P_a$, the previously selected point from the preceding bucket, and $P_b$, the averaged representative point (mean index and mean intensity) of the next bucket. Here $i_a, i_c, i_b$ denote the channel indices and $I_a, I_c, I_b$ the corresponding intensities of $P_a, P_c$, and $P_b$ respectively. The signed triangle area is computed according to Equation~\eqref{eq:lttb}:
\begin{equation}
    \label{eq:lttb}
    A(P_a, P_c, P_b) = \frac{1}{2} |i_a(I_c - I_b) + i_c(I_b - I_a) + i_b(I_a - I_c)|
\end{equation}
The candidate $P_c$ within bucket $j$ maximising $|A(P_a, P_c, P_b)|$ is selected, and the anchor $P_a$ updates to the newly selected point before evaluating the next bucket. However, geometric downsampling algorithms like LTTB do not guarantee peak preservation.

To prevent the loss of narrow diagnostic emission lines, LunarAtlas implements a Targeted Peak Preservation (Peak-Union) step. The Peak-Union procedure operates as follows:
\begin{enumerate}
  \item Retrieve the set of curated NIST reference wavelengths $P_{\text{NIST}}$ corresponding to target elements under investigation (e.g. Ca II at 393.366~nm, Mg II at 279.553~nm, etc.).
  \item Run standard LTTB on the full spectrum to produce the downsampled index set $S_{\text{LTTB}}$.
  \item Compute $S_{\text{final}} = S_{\text{LTTB}} \cup \{i_p : p \in P_{\text{NIST}}\}$, where $i_p$ is the nearest channel index to each reference wavelength.
  \item Sort $S_{\text{final}}$ by wavelength index.
\end{enumerate}
The output size is therefore $|S_{\text{final}}| \leq |S_{\text{LTTB}}| + |P_{\text{NIST}}|$, with equality when no target peaks fall on LTTB-selected indices.

The peak retention rate measures the percentage of originally detected spectral peaks that are preserved in the downsampled subset (either exactly or at adjacent channels).

\subsubsection{NIST Reference Matching Criterion}
LunarAtlas integrates reference wavelengths from the NIST Atomic Spectra Database (ASD) \cite{nist_asd} to support interactive spectral exploration and candidate wavelength matching. A peak at $\lambda_{\mathrm{obs}}$ is classified as a candidate match if and only if: (i) it is identified as a local maximum under a $3\sigma$ prominence threshold, where $\sigma$ is the population standard deviation ($\mathit{ddof}=0$) computed over the 560--580~nm continuum window of the same spectrum; and (ii) $|\lambda_{\mathrm{obs}} - \lambda_{\mathrm{NIST}}| \leq 0.5$~nm \cite{sridhar2024libs} for the nearest NIST reference line. Both conditions must be satisfied simultaneously.

The $\pm0.5$~nm tolerance was adopted to accommodate the spectral resolution of the Level-1 dataset and minor wavelength calibration uncertainties. The LE-LIBS instrument has a stated optical resolution of $\leq 1$~nm; the adopted tolerance therefore corresponds to half a resolution element --- a relatively permissive criterion. Channels with intensity above $3\sigma$ but sub-threshold prominence --- residing within broad elevated continuum regions rather than constituting resolved emission features --- do not satisfy condition~(i) and are excluded from programmatic matching. They are rendered as reference overlays in the interface but are explicitly distinguished from candidate peak matches.

Confirmation of elemental composition requires corroborating evidence beyond wavelength coincidence, including line ratio analysis, cross-validation with additional transitions of the same species, plasma temperature constraints, and independent spectroscopic validation. Matches reported under this criterion therefore support exploratory candidate identification only.

\begin{figure*}[!t]
    \centering
    \includegraphics[width=0.9\textwidth]{publication_figures/fig2_reshape_schematic}
    \caption{L1 Wide-Format Column Structure to Normalised Long-Form Schema Ingestion.}
    \label{fig:reshape_schematic}
\end{figure*}


\section{Illustrative examples}

A typical workflow involves cloning the repository, installing dependencies, and executing the multi-stage ingestion pipeline. The live deployment is accessible at \url{https://www.lunaratlas.in} and the versioned REST API is served at \url{https://api.lunaratlas.in}. These live services are hosted for interactive demonstration purposes; the primary, citable, and reproducible artifact of this work is the versioned codebase archived at Zenodo \cite{lunaratlas_zenodo}. The following CLI commands illustrate local environment setup, pipeline execution, and server startup:

\begin{lstlisting}[language=bash, caption=Environment setup, pipeline execution, and server boot commands]
# Install server and client dependencies
pip install -r core/server/requirements.txt
cd core/client && npm install && cd ../..

# Execute the 6-stage data processing and ingestion pipeline
cd Pipeline && ./run_pipeline.sh

# Start the FastAPI backend API server
cd core/server && uvicorn app.main:app --host 0.0.0.0 --port 8000

# Start the React client developer server (in a separate shell)
cd core/client && npm run dev
\end{lstlisting}

Once the services are active, the REST API can be queried programmatically. The backend server exposes several endpoints: (i) \texttt{/api/v1/spectrum} which retrieves clean spectral curves with adaptive LTTB+Peaks downsampling and optional denoising filters (ALS, Savitzky-Golay); (ii) public index queries such as \texttt{/api/v1/public/observations} and \texttt{/api/v1/public/measurements} which expose PDS4 metadata and physical parameters (laser energy, pulses) filtered by date or identifier; and (iii) status monitors such as \texttt{/api/v1/health}.

The following script queries the public API to retrieve a peak-locked downsampled spectrum for a processed measurement:

\begin{lstlisting}[language=Python, caption=Python client request for downsampled spectra]
import requests

url = "https://api.lunaratlas.in/api/v1/spectrum"
params = {
    "measurement_id": "1",  # clean Measurement ID
    "lambda_min": 240.0,
    "lambda_max": 450.0,
    "zoom_level": 2,
    "als": True,
    "savgol": True
}

response = requests.get(url, params=params)
data = response.json()

# Process spectrum points (DownsampledBucket or SpectralDataPoint objects)
for pt in data["data"][:5]:
    if "wavelength_nm" in pt:
        print(f"Wavelength: {pt['wavelength_nm']} nm, Cleaned: {pt['intensity']} a.u.")
    else:
        print(f"Bucket: {pt['lambda_min']} to {pt['lambda_max']} nm")
\end{lstlisting}

Figure~\ref{fig:ui_console} illustrates the interactive visualisation console supported by the LTTB+Peaks algorithm. The interface shows the spectral view with NIST reference line overlays for Measurement ID~3, the highest-signal acquisition (mean 417~a.u.). The visualization dynamically updates using LTTB+Peaks downsampling, transitioning to raw un-decimated mode when the viewport is zoomed in and the channel count falls below the target threshold.

\begin{figure*}[!t]
    \centering
    \includegraphics[width=0.9\textwidth]{publication_figures/fig5_nist_overlay}
    \caption{UI visualisation console showing NIST reference line overlays on the cleaned LE-LIBS spectrum for Measurement ID~3. Labelled species include confirmed peak matches Mg~II (279.55/280.27~nm) and Ca~II (393.37/396.85~nm); the sub-threshold transition Si~I (288.16~nm) is displayed as a visual overlay only.}
    \label{fig:ui_console}
\end{figure*}

\section{Impact}

\textit{Scope note.} All quantitative results in this section constitute a single worked example performed on the Chandrayaan-3 LE-LIBS Level-1 product \texttt{ch3\_lib\_002\_20230825T104221\_00\_l1}. The intent is to demonstrate that the infrastructure produces physically plausible and internally consistent outputs on a representative input file. Systematic generalisation across the full PRADAN archive has not been tested; the numerical values should therefore not be interpreted as population-level benchmarks.

\subsection{Cleaning Pipeline Validation}

All experimental evaluations and benchmarks in this section were conducted on the Chandrayaan-3 Pragyan Rover LE-LIBS Level-1 (L1) calibrated product \texttt{ch3\_lib\_002\_20230825T104221\_00\_l1} (corresponding to raw file \texttt{ch3\_lib\_002\_20230825T104221\_00\_l1.csv} containing 8 interleaved acquisitions, representing 4 background and 4 plasma measurements).

The LunarAtlas cleaning stage applies two preprocessing operations: (i) subtraction of the most recently acquired background spectrum preceding each plasma shot, and (ii) non-negativity enforcement through zero-clamping. Background subtraction is a standard spectroscopic preprocessing technique used to remove instrumental baseline and continuum contributions prior to spectral analysis \cite{cremers2006libs, chen2023automatic}.

The cleaning pipeline was evaluated on a 20~nm validation window (560--580~nm), a spectral region free of major expected lunar regolith emission lines (such as strong transitions of Fe, Ti, Ca, Al, or Mg) and representative of the instrumental continuum baseline, as validated against the NIST ASD reference database \cite{nist_asd} and standard LIBS regolith references \cite{cremers2006libs, sridhar2024libs}. In this window, the raw L1 spectrum for Measurement ID~3 had a mean intensity of $\approx1{,}387$~a.u. and a standard deviation of $\approx54$~a.u. In contrast, the cleaned spectrum had a mean intensity of $\approx444$~a.u. and a standard deviation of $\approx71$~a.u. This yields a baseline suppression factor, defined in Equation~\eqref{eq:suppression}:
\begin{equation}
    \label{eq:suppression}
    \text{Baseline suppression} = \bar{I}_{\mathrm{raw}} / \bar{I}_{\mathrm{clean}} \approx 3.12
\end{equation}
The $3.12\times$ suppression demonstrates the removal of the background continuum (Figure~\ref{fig:pre_post_overlay}). For reference, background-to-plasma intensity ratios of 2--5${\times}$ are routinely reported in laboratory LIBS measurements of geological silicate targets \cite{cremers2006libs}, and are consistent with the continuum emission characteristics of a pulsed Yb:Er:Glass laser at the operating fluence of the Pragyan instrument \cite{sridhar2024libs}. The observed $3.12{\times}$ factor therefore falls within the physically expected range for this instrument and target class, and is consistent with successful background suppression rather than anomalous behaviour.

As an independent correctness check, the standard deviation ratio between cleaned and raw spectra in the validation window was 1.31, compared to the theoretical $\sqrt{2} \approx 1.41$ expected from the subtraction of two independent noise signals \cite{bevington2003data}. The 7.7\% discrepancy below the theoretical limit is expected due to minor positive correlation in sequential CCD readouts and the narrow 20~nm evaluation window --- both of which reduce the effective noise independence assumed in the theoretical derivation. This agreement constitutes a quantitative verification that the background subtraction is operating correctly. Prior to clamping, approximately 0.5\% of all spectral channels for Measurement ID~3 contained negative intensity values due to statistical fluctuations. Since detector response intensity values (a.u.) are physically non-negative quantities, all negative values were clamped to zero. Table~\ref{table:ablation} summarises the pipeline metrics at each stage of the cleaning process.

\begin{figure}[!htbp]
    \centering
    \includegraphics[width=\columnwidth]{publication_figures/fig3_pre_post_overlay}
    \caption{Pre- vs.\ post-subtraction spectrum overlay in the 560--580~nm validation window for Measurement ID~3.}
    \label{fig:pre_post_overlay}
\end{figure}

\begin{table}[!htbp]
\centering
\footnotesize
\setlength{\tabcolsep}{4pt}
\begin{tabular}{@{}p{2.8cm}p{2.8cm}c@{}}
\toprule
\textbf{Stage} & \textbf{Metric} & \textbf{Value} \\
\midrule
Raw L1 Input (ID~3) & Window Mean Intensity & 1,387~a.u. \\
Paired BG Subtraction & Suppression Ratio & 3.12$\times$ \\
Prior to Clamping & Negative Channels & 0.5\% \\
Zero-Clamping & Negative Channels & 0.0\% \\
\bottomrule
\end{tabular}
\caption{Ingestion pipeline stage-by-stage validation metrics in the 560--580~nm window for Measurement ID~3.}
\label{table:ablation}
\end{table}


\subsection{Inter-Measurement Variability}

LunarAtlas preserves individual Measurement IDs rather than averaging multiple plasma shots into a single spectrum. Table~\ref{table:variability} demonstrates substantial shot-to-shot intensity variability within the same observation sequence. For example, Measurement ID~3 exhibited a mean intensity of 416.9~a.u. with only 0.5\% clamped channels, whereas Measurement ID~2 exhibited a mean intensity of 9.4~a.u. with 77.4\% clamped channels, indicating a low-signal or potentially failed plasma event. If these measurements were averaged together, such low-signal acquisitions would be partially obscured. Preserving measurement-level granularity therefore enables analysis of shot-to-shot variability, identification of anomalous plasma events, and quality assessment of individual acquisitions.

\begin{table}[!htbp]
\centering
\footnotesize
\setlength{\tabcolsep}{4pt}
\begin{tabular}{@{}ccccc@{}}
\toprule
\textbf{ID} & \textbf{Mean} & \textbf{Std} & \textbf{Max} & \textbf{Clamped} \\
 & \textbf{(a.u.)} & \textbf{(a.u.)} & \textbf{(a.u.)} & \textbf{(\%)} \\
\midrule
1 & 61.0  & 56.9 & 321 & 23.0 \\
2 & 9.4   & 23.8 & 261 & 77.4 \\
3 & 416.9 & 95.0 & 684 & 0.5  \\
4 & 54.1  & 62.4 & 637 & 28.7 \\
\bottomrule
\end{tabular}
\caption{Cleaned intensity statistics and zero-clamping fraction per Measurement ID (all 2,094 channels). ID~2 (77.4\% clamped) is a suspected low-signal or failed plasma event.}
\label{table:variability}
\end{table}

The observed shot-to-shot variability (mean intensities ranging from 9.4 to 416.9~a.u.) may reflect genuine plasma physics --- such as surface heterogeneity, laser coupling efficiency variation, or ablation crater evolution --- as well as instrumental artefacts including detector read noise and laser pulse energy fluctuations. With only four acquisitions in this observation sequence, distinguishing these contributions is not possible; the variability is therefore reported descriptively rather than interpreted causally. Multi-session analysis across the full PRADAN archive would be required for a statistically powered characterisation.

\subsection{Downsampling and Peak Preservation}

Evaluation of the downsampling algorithms was performed on valid plasma measurements (IDs 1, 3, and 4) which exhibit clear emission features. Measurements ID~1, ID~3, and ID~4 were selected because they contain sufficient spectral peaks for quantitative evaluation, whereas low-signal acquisitions (e.g., ID~2) do not permit meaningful peak-preservation analysis. ID~2 (77.4\% clamped channels, exceeding the 50\% quality-warning threshold) was excluded for this reason.

Peak detection was performed using \texttt{scipy.signal.find\_peaks} with a prominence threshold of $3\sigma_{\text{noise}}$ (where $\sigma_{\text{noise}}$ is the population standard deviation, $\mathit{ddof}{=}0$, in the 560--580~nm continuum window) and a minimum inter-peak distance of 5 channels. This yielded 76 local maxima in Measurement ID~3 and 91 in Measurement ID~4, with prominence thresholds of $3\sigma = 212.28$~a.u. and $140.65$~a.u. respectively.

The LTTB+Peaks downsampling algorithm was compared against other decimation algorithms at 10\% data density, as summarized in Table~\ref{table:downsampling}. The Local Ret.\ (\%) and NIST Ret.\ (\%) columns report the fraction of detected local maxima and NIST-matched peaks preserved by each algorithm, respectively. RMSE is computed relative to the full-resolution 2,094-point spectrum via linear interpolation of the downsampled points back to full resolution. RMSE is computed via linear interpolation of downsampled points back to full resolution; this interpolation introduces artificial smoothing at narrow emission lines, which modestly inflates RMSE for peak-preserving algorithms. RMSE is therefore reported as a measure of global spectral envelope fidelity, complementary to — not competing with — peak retention rate.

\begin{table*}[!htbp]
\centering
\footnotesize
\setlength{\tabcolsep}{6pt}
\begin{tabular}{@{}lcrrr@{}}
\toprule
\textbf{Algorithm} & \textbf{Meas. ID} & \textbf{Local ret. (\%)} & \textbf{NIST ret. (\%)} & \textbf{RMSE (a.u.)} \\
\midrule
            & 1 & 4.8  & 0.0   & 62.83  \\
Uniform     & 3 & 11.9 & 0.0   & 78.87  \\
            & 4 & 12.3 & 0.0   & 62.95  \\
\midrule
            & 1 & 100.0 & 100.0 & 106.24 \\
Max-binning & 3 & 98.3  & 100.0 & 118.38 \\
            & 4 & 87.7  & 100.0 & 102.05 \\
\midrule
            & 1 & 71.4 & 100.0 & 71.79  \\
Std. LTTB   & 3 & 47.5 & 57.1  & 85.30  \\
            & 4 & 56.2 & 50.0  & 68.17  \\
\midrule
\textbf{LTTB+Peaks} & 1 & \textbf{100.0} & \textbf{100.0} & \textbf{73.42}  \\
(proposed)          & 3 & \textbf{100.0} & \textbf{100.0} & \textbf{87.49}  \\
                    & 4 & \textbf{100.0} & \textbf{100.0} & \textbf{73.23}  \\
\bottomrule
\end{tabular}
\caption{Downsampling performance comparison at 10\% density (209 of 2,094~points) across all valid plasma measurements (Measurement IDs 1, 3, and 4). Local ret.\ is the fraction of detected local maxima preserved. NIST ret.\ is the fraction of NIST reference lines preserved. RMSE is computed against the full-resolution spectrum via linear interpolation.}
\label{table:downsampling}
\end{table*}

RMSE quantifies global spectral shape fidelity---how well the downsampled envelope approximates the full-resolution spectrum. Peak retention quantifies feature-level fidelity---whether diagnostically important emission features survive decimation. Both metrics are complementary: an algorithm may achieve low RMSE while discarding narrow peaks, or preserve all peaks while distorting the spectral envelope. Uniform decimation achieves the lowest latency but retains only 13.2\% of detected peaks, as spectral emission lines are not uniformly distributed across the wavelength axis. Max-binning achieves high peak retention (96.1\%) by selecting the maximum value within each bin, but incurs a higher RMSE (121.24~a.u.) because bin maxima selection distorts the spectral envelope between peaks. While standard LTTB achieves a lower RMSE (90.10~a.u.), it only retains 52.6\% of local maxima. The addition of the Peak-Union step in LTTB+Peaks achieves 100\% retention of detected peaks with a small RMSE penalty (92.24~a.u.). The RMSE difference between standard LTTB (90.10~a.u.) and LTTB+Peaks (92.24~a.u.) is 2.14~a.u.\ (2.4\% relative), attributable in part to interpolation artefacts at preserved peak positions; this penalty is small relative to the 212.28~a.u.\ prominence threshold used for peak detection. To confirm that LTTB+Peaks retention is not specific to a single acquisition, the evaluation was repeated on all valid plasma measurements (IDs 1, 3, and 4; ID 2 excluded, 77.4\% clamped channels). LTTB+Peaks achieved 100\% local maxima and NIST peak retention across all three spectra and all tested densities (Tables 5 and 6), with RMSE penalties relative to standard LTTB of 1.63–5.06 a.u. (2.3–7.3\% relative), confirming the result is not artefactual to a single spectrum.

To validate peak preservation from both an engineering perspective (all detected peaks) and a spectroscopy-oriented perspective (NIST reference lines), Table~\ref{table:downsampling} reports both local maxima and NIST-matched peak retention rates. A matching tolerance of $\pm0.5$~nm was adopted to accommodate the spectral resolution of the Level-1 dataset and minor wavelength calibration uncertainties. Out of 76 detected local maxima in Measurement ID~3, 7 peaks matched NIST reference wavelengths within this tolerance: Fe~I (274.86~nm), Ca~II (393.33~nm), Ca~II (396.94~nm), Fe~I (496.01~nm), Mg~I (517.25~nm), Na~I (589.15~nm), and O~I (777.67~nm). The standard LTTB algorithm achieved 57.1\% NIST retention (preserving 4 of these 7 lines), whereas LTTB+Peaks preserved all 7 NIST-matched lines (100.0\% retention) by explicitly appending detected peak indices to the downsampled set. The small number of NIST-matched peaks (7 of 76 detected) reflects the stringent dual requirement — both the 3\sigma prominence threshold and the \pm0.5 nm matching tolerance must be satisfied simultaneously. The Peak-Union step guarantees 100\% retention of these indices by construction; the informative claim is therefore not that 7 peaks are preserved, but that the RMSE cost of doing so is 2.4\% relative to standard LTTB, and that this guarantee holds across all tested densities (Table 6) and all valid acquisitions (Table 5).

To demonstrate that this retention holds across different target compression levels, Table~\ref{table:density_sweep} displays peak retention as a function of downsampling density (from 5\% to 50\% points). LTTB+Peaks consistently achieves 100.0\% retention across all tested densities. Standard LTTB retention degrades to 39.5\% at 5\% density. Because the downsampling algorithm explicitly appends the original peak indices, the retained peaks exhibit zero measured position and amplitude error under all decimation ratios.

\begin{table*}[!htbp]
\centering
\footnotesize
\setlength{\tabcolsep}{8pt}
\begin{tabular}{@{}lcrrrrr@{}}
\toprule
\textbf{Algorithm} & \textbf{Meas. ID} & \textbf{5\%} & \textbf{10\%} & \textbf{15\%} & \textbf{20\%} & \textbf{50\%} \\
\midrule
           & 1 & \textbf{100.0} & \textbf{100.0} & \textbf{100.0} & \textbf{100.0} & \textbf{100.0} \\
LTTB+Peaks & 3 & \textbf{100.0} & \textbf{100.0} & \textbf{100.0} & \textbf{100.0} & \textbf{100.0} \\
           & 4 & \textbf{100.0} & \textbf{100.0} & \textbf{100.0} & \textbf{100.0} & \textbf{100.0} \\
\midrule
           & 1 & 50.0 & 71.4 & 66.7 & 71.4 & 90.5 \\
Std. LTTB  & 3 & 39.0 & 47.5 & 62.7 & 57.6 & 94.9 \\
           & 4 & 39.7 & 56.2 & 67.1 & 72.6 & 95.9 \\
\bottomrule
\end{tabular}
\caption{Local peak retention (\%) as a function of downsampling density (fraction of 2,094 original channels) for LTTB+Peaks vs. standard LTTB across all valid plasma measurements (Measurement IDs 1, 3, and 4; ID 2 is excluded due to the quality-warning threshold).}
\label{table:density_sweep}
\end{table*}

\begin{figure}[!htbp]
    \centering
    \includegraphics[width=\columnwidth]{publication_figures/fig4_peak_comparison}
    \caption{Comparison of peak preservation by standard LTTB and LTTB+Peaks.}
    \label{fig:peak_preservation}
\end{figure}

\subsection{NIST-Assisted Spectral Annotation --- Applied Results}

The matching criterion is defined in Section 2.2.3. Applying both conditions across the three valid plasma measurements, candidate matches were identified in all spectra (Table~\ref{table:elements}): 8 of 59 detected peaks in Measurement ID~3 ($13.6\%$), 2 of 42 in Measurement ID~1 ($4.8\%$), and 5 of 73 in Measurement ID~4 ($6.8\%$). The varying match rate reflects each spectrum's noise floor --- Measurement ID~3's higher signal-to-noise ratio ($\sigma = 70.76$~a.u.) yields a more permissive $3\sigma$ threshold relative to genuine emission features than the lower-signal ID~1 ($\sigma = 59.17$~a.u.) or ID~4 ($\sigma = 46.88$~a.u.).

Five additional NIST lines --- Mg~I (285.21~nm), Si~I (288.16~nm), Al~I (309.27~nm), Fe~I (404.58~nm), and Ca~I (422.67~nm) --- have intensities of 387--572~a.u. but sub-threshold prominence and are therefore visual overlays only. The H~I Balmer-$\alpha$ line (NIST: 656.28~nm) corresponds to a detected peak at 657.09~nm (prominence 221.0~a.u.) but its offset of 0.81~nm exceeds the $\pm0.5$~nm criterion; it is retained as a flagged near-match with the offset explicitly displayed.

\begin{table*}[!htbp]
\centering
\footnotesize
\setlength{\tabcolsep}{6pt}
\begin{tabular}{@{}lccccr@{}}
\toprule
\textbf{Species} & \textbf{$\lambda_{\mathrm{NIST}}$ (nm)} & \textbf{$\lambda_{\mathrm{obs}}$ (nm)} & \textbf{Int. (a.u.)} & \textbf{Prom. (a.u.)} & \textbf{$\Delta\lambda$ (nm)} \\
\midrule
\multicolumn{6}{l}{\textbf{ID 1} \quad $\sigma = 59.17$~a.u. $\cdot$ $3\sigma = 177.51$~a.u. $\cdot$ 42 peaks detected $\cdot$ 2 matched ($4.8\%$)} \\
\midrule
Ca~II & 396.85 & 396.94 & 215 & 215.0 & 0.09 \\
O~I   & 777.19 & 777.05 & 213 & 213.0 & 0.14 \\
\midrule
\multicolumn{6}{l}{\textbf{ID 3} \quad $\sigma = 70.76$~a.u. $\cdot$ $3\sigma = 212.28$~a.u. $\cdot$ 59 peaks detected $\cdot$ 8 matched ($13.6\%$)} \\
\midrule
Fe~I  & 274.95 & 274.86 & 525 & 242.0 & 0.09 \\
Fe~I  & 385.99 & 386.47 & 462 & 238.0 & 0.48 \\
Ca~II & 393.37 & 393.33 & 491 & 247.0 & 0.04 \\
Ca~II & 396.85 & 396.94 & 495 & 217.0 & 0.09 \\
Fe~I  & 495.76 & 496.01 & 583 & 244.0 & 0.25 \\
Mg~I  & 517.27 & 517.25 & 640 & 401.0 & 0.02 \\
Na~I  & 588.99 & 589.15 & 561 & 254.0 & 0.16 \\
O~I   & 777.19 & 777.67 & 568 & 222.0 & 0.48 \\
\midrule
\multicolumn{6}{l}{\textbf{ID 4} \quad $\sigma = 46.88$~a.u. $\cdot$ $3\sigma = 140.65$~a.u. $\cdot$ 73 peaks detected $\cdot$ 5 matched ($6.8\%$)} \\
\midrule
Fe~I  & 274.95 & 274.86 & 191 & 173.0 & 0.09 \\
Mg~I  & 285.21 & 284.90 & 251 & 173.0 & 0.31 \\
Si~I  & 288.16 & 288.25 & 178 & 178.0 & 0.09 \\
Mg~I  & 517.27 & 516.90 & 213 & 213.0 & 0.37 \\
O~I   & 777.19 & 777.05 & 194 & 194.0 & 0.14 \\
\bottomrule
\end{tabular}
\caption{Candidate wavelength matches --- NIST ASD vs.\ detected local maxima across valid plasma measurements (Measurement IDs 1, 3, and 4; ID 2 is excluded due to the quality-warning threshold). Only peaks satisfying both the $3\sigma$ prominence condition and the $\pm0.5$~nm tolerance \cite{sridhar2024libs} are included. Int.: peak intensity; Prom.: peak prominence. Matches support exploratory analysis only and do not constitute definitive elemental identifications.}
\label{table:elements}
\end{table*}

\textit{Interpretation.} Confirmation of elemental composition requires corroborating evidence beyond wavelength coincidence, including line ratio analysis, cross-validation with additional transitions of the same species, plasma temperature constraints, and independent spectroscopic validation.

\subsection{API Response Performance}

All benchmarks were conducted on a single-node workstation (Intel Core i7-11850H, 16~GB RAM, NVMe SSD) running Ubuntu 22.04, Python 3.11, and PostgreSQL~16.3 (or SQLite co-located for local data-path isolation) under no-load, single-connection conditions.

To characterize the baseline processing budget of the software stack (FastAPI routing, schema validation, LTTB+Peaks downsampling, and JSON serialization) independent of geographical WAN transit overhead, we executed local benchmarks targeting the spectrum retrieval endpoint (\texttt{/api/v1/public/spectra/\{measurement\_id\}}) on port 8020. With local database co-location to isolate the processing stack, we performed 1,000 measured runs (following 20 discarded warm-up runs) for all three valid plasma acquisitions:
\begin{itemize}
    \item \textbf{Measurement ID 1}: Mean end-to-end latency was $34.18 \pm 27.85$~ms (median: $23.32$~ms; 95th percentile: $99.42$~ms; 99th percentile: $128.27$~ms). The $3.26 \pm 0.25$~ms LTTB+Peaks downsampling step contributed approximately $9.54\%$ of this total.
    \item \textbf{Measurement ID 3}: Mean end-to-end latency was $35.65 \pm 29.06$~ms (median: $24.60$~ms; 95th percentile: $104.54$~ms; 99th percentile: $133.80$~ms). The downsampling step contributed approximately $9.14\%$ of this total.
    \item \textbf{Measurement ID 4}: Mean end-to-end latency was $36.68 \pm 30.30$~ms (median: $24.97$~ms; 95th percentile: $107.94$~ms; 99th percentile: $144.79$~ms). The downsampling step contributed approximately $8.89\%$ of this total.
\end{itemize}
These unimodal, right-skewed latency profiles confirm that the LTTB+Peaks downsampling algorithm (characterized in Table~\ref{table:downsampling}) contributes minimally (under $10\%$) to the local API processing budget, and that the dominant local latency components are database query execution and data serialization.

For comparison, we benchmarked the same endpoint on the live production server (\url{https://api.lunaratlas.in}) located in a remote datacenter, which incorporates internet transit. Under single-connection conditions, the mean production round-trip latency to fetch the complete spectrum was $983.62$~ms. Under concurrent-load testing on the production server, the mean round-trip latency scaled to $1{,}049.03$~ms for 5 concurrent connections, $1{,}053.79$~ms for 10 concurrent connections, and $1{,}377.61$~ms for 20 concurrent connections. The sub-linear latency growth demonstrates that the production environment benefits from connection pooling and Redis caching, establishing the system's viability for collaborative concurrent scientific use.


\begin{table}[!htbp]
\centering
\footnotesize
\setlength{\tabcolsep}{4pt}
\begin{tabular}{@{}p{4.2cm}c@{}}
\toprule
\textbf{Metric} & \textbf{Value} \\
\midrule
Max.\ API Latency & $<$500~ms \\
DB Query Speedup (index vs.\ seq.\ scan) & $\approx 85{\times}$ \\
Overview Payload Size & 10--400~kB \\
\bottomrule
\end{tabular}
\caption{System performance metrics (Intel Core i7-11850H, 16~GB RAM, NVMe SSD, Ubuntu 22.04, Python~3.11, PostgreSQL~16.3, single-connection, no-load).}
\label{table:performance}
\end{table}

\begin{figure*}[!t]
    \centering
    \includegraphics[width=0.95\textwidth]{publication_figures/fig6_adaptive_zoom}
    \caption{LTTB+Peaks downsampling applied to the cleaned Chandrayaan-3 spectrum (Measurement ID~3) at four zoom levels ($k=1, 2, 3,$ and $4$), showing full-resolution references (light grey), downsampled curves (blue), and preserved peaks (red markers).}
    \label{fig:adaptive_zoom}
\end{figure*}

\section{Conclusions}

LunarAtlas provides an open-source, reproducible data processing and database infrastructure for Chandrayaan-3 LE-LIBS Level-1 observations. The pipeline implements a deterministic processing chain that resolves the challenges of tabular PDS4 parsing, interleaved acquisition pairing, and physically motivated background cleaning---validated through quantitative evaluation of cleaning, peak preservation, and API performance. The peak-locked LTTB+Peaks downsampling algorithm achieved 100\% peak retention across all evaluated spectra and density levels (5--50\%, evaluated over the two valid plasma acquisitions in the reference product) during visualisation while maintaining sub-500~ms response times on commodity hardware. Pipeline correctness is independently supported by the observed noise propagation factor (1.31 vs.\ theoretical 1.41), consistent with the expected behaviour of paired background subtraction on correlated CCD readouts. By establishing an auditable, versioned reference archive for Chandrayaan-3 LIBS data, LunarAtlas supports quantitative lunar science and offers a template architecture applicable to future planetary LIBS mission archives. The system is publicly accessible at \url{https://www.lunaratlas.in}, with the REST API at \url{https://api.lunaratlas.in}.

\paragraph{Limitations.} Benchmarks and evaluations in Section~4 were performed on a single representative reference product containing four cleaned spectral measurements, and are therefore presented as a worked example rather than a generalised benchmark; statistical claims about inter-acquisition variability should therefore be interpreted as software validation rather than a comprehensive analysis of lunar LIBS measurements. No Level-2 calibration is implemented --- LunarAtlas provides analysis-ready L1 records as input to such workflows. The system has not been benchmarked under concurrent multi-user load.

\paragraph{Future work.} The most scientifically significant open direction is Level-2 calibration: translating cleaned L1 intensity records into quantitative elemental abundance estimates requires plasma temperature modelling, matrix effect correction, and validation against certified geological reference standards. On the infrastructure side, horizontal scaling for concurrent query support remains an open engineering task as community adoption grows. More broadly, LunarAtlas is intended as a template architecture for future planetary LIBS archives; its ingestion and cleaning abstractions are designed to generalise to forthcoming missions where in-situ LIBS will remain a primary geochemical tool.

\section*{Author Contributions}
\textbf{Lovekesh Anand:} Conceptualisation, Methodology, Software, Validation, Formal Analysis, Data Curation, Writing -- Original Draft, Visualisation.
\textbf{Dua Saeed:} Writing -- Review \& Editing, Supervision.

\section*{Declaration of Competing Interests}
The authors declare that they have no known competing financial interests or personal relationships that could have appeared to influence the work reported in this paper.

\section*{Funding}
This research received no specific grant from any funding agency in the public, commercial, or not-for-profit sectors.

\section*{Acknowledgements}
The authors thank the Indian Space Research Organisation (ISRO) for making the Chandrayaan-3 LE-LIBS Level-1 data public through the PRADAN data portal. Wavelength reference data were accessed from the NIST Atomic Spectra Database.

\section*{Software Availability}
\begin{itemize}
    \item \textbf{Developer / Authors:} Lovekesh Anand, Dua Saeed
    \item \textbf{Contact address:} Maharshi Dayanand University, Rohtak, Haryana, India
    \item \textbf{E-mail:} \href{mailto:lovekeshanand6@gmail.com}{lovekeshanand6@gmail.com}, \href{mailto:23f1000825@ds.study.iitm.ac.in}{23f1000825@ds.study.iitm.ac.in}
    \item \textbf{Software repository:} \url{https://github.com/LovekeshAnand/LunarAtlas}
    \item \textbf{Permanent archive identifier (Zenodo DOI):} \url{https://doi.org/10.5281/zenodo.20771541} \cite{lunaratlas_zenodo}
    \item \textbf{License:} MIT License
    \item \textbf{First released:} 2026
\end{itemize}

\section*{Data Availability}
The Chandrayaan-3 LE-LIBS Level-1 data used in this study are publicly available from the ISRO Planetary Data Archive (PRADAN) at \url{https://pradan.issdc.gov.in}. The live system is accessible at \url{https://www.lunaratlas.in} and the versioned REST API at \url{https://api.lunaratlas.in} (provided for interactive demonstration purposes; the primary citable codebase is the Zenodo repository listed above).

\bibliographystyle{elsarticle-num}
\bibliography{paper}

\end{document}
\endinput